% GENERATED FILE. Edit canonical lessons, then run npm run generate:books.
% canonical-source-hash: fnv1a64:6c6e9b4f4f2ddf78
% canonical-lessons: RU-C478-upryamyy, RU-C478-upryamstvo, RU-C478-kapriznyy, RU-C478-kapriznichat, RU-C478-vorchat

\chapter{Stubborn and Fussy}
\label{ch:ru-stubborn-and-fussy}
\hlchaptermodality{478}

\begin{chapteropening}
\textbf{By the end of this chapter:} \emph{I can describe stubborn, fussy and grumpy people.}

Complete the last lesson of chapter 478: I can describe stubborn, fussy and grumpy people.
\end{chapteropening}

\section[\texorpdfstring{\ru{упрямый} (upryámyy)}{упрямый (upryámyy)}]{\ru{упрямый} (upryámyy) — stubborn}

\label{lesson:RU-C478-upryamyy}



\begin{quote}
Before the new one: say the Russian for lonely, then the Russian for to regret.
\end{quote}

\subsection*{You'll want to know: \ru{упрямый}}
\textbf{\ru{упрямый}} (\emph{upryámyy}) — \textquotedblleft{}stubborn\textquotedblright{}.

\begin{quote}
\ru{Упрямый} \ru{ребёнок} \ru{не} \ru{хотел} \ru{есть} \ru{суп}.
\end{quote}

\begin{quote}
\emph{Upryámyy rebyónok ne khotél yest sup.}
\end{quote}

\begin{quote}
The stubborn child wouldn't eat the soup.
\end{quote}

Short forms: \ru{упрям}, \ru{упряма}.

\begin{cousinweb}
EtymDB lists \emph{\ru{упрямый}} without a source; it is usually linked to \emph{\ru{упереться}} 'to dig one's heels in', but the origin is not settled.
\end{cousinweb}

\begin{grammarlens}[title={What you've built}]
One more word for telling, explaining and describing at B1.
\end{grammarlens}

\subsection*{Your turn}
\begin{itemize}
\raggedright
  \item \emph{Say it:} \emph{upryámyy}
  \item \emph{Say it:} \emph{upryámyy}, once more
  \item \emph{Say it:} a sentence of your own with \emph{upryámyy}
  \item \emph{Recall:} say the Russian for lonely, then the Russian for to regret, then say \emph{upryámyy} again
\end{itemize}

\begin{tcolorbox}[breakable,colback=teal!4,colframe=teal!35!black,title={Before you move on}]
What does \ru{упрямый} mean? (\textquotedblleft{}Stubborn\textquotedblright{}.) Say it once more.
\end{tcolorbox}
\section[\texorpdfstring{\ru{упрямство} (upryámstvo)}{упрямство (upryámstvo)}]{\ru{упрямство} (upryámstvo) — stubbornness}

\label{lesson:RU-C478-upryamstvo}



\begin{quote}
Before the new one: say the Russian for to regret, then the Russian for stubborn.
\end{quote}

\subsection*{You'll want to know: \ru{упрямство}}
\textbf{\ru{упрямство}} (\emph{upryámstvo}) — \textquotedblleft{}stubbornness\textquotedblright{}.

\begin{quote}
\ru{Его} \ru{упрямство} \ru{нас} \ru{удивило}.
\end{quote}

\begin{quote}
\emph{Yegó upryámstvo nas udivílo.}
\end{quote}

\begin{quote}
His stubbornness surprised us.
\end{quote}

Neuter.

\begin{cousinweb}
\emph{\ru{упрямство}} is built on \emph{\ru{упрямый}} 'stubborn' with \emph{-\ru{ство}}; EtymDB has no entry for the noun, and the origin of \emph{\ru{упрямый}} is not settled.
\end{cousinweb}

\begin{grammarlens}[title={What you've built}]
One more word for telling, explaining and describing at B1.
\end{grammarlens}

\subsection*{Your turn}
\begin{itemize}
\raggedright
  \item \emph{Say it:} \emph{upryámstvo}
  \item \emph{Say it:} \emph{upryámstvo}, once more
  \item \emph{Say it:} a sentence of your own with \emph{upryámstvo}
  \item \emph{Recall:} say the Russian for to regret, then the Russian for stubborn, then say \emph{upryámstvo} again
\end{itemize}

\begin{tcolorbox}[breakable,colback=teal!4,colframe=teal!35!black,title={Before you move on}]
What does \ru{упрямство} mean? (\textquotedblleft{}Stubbornness\textquotedblright{}.) Say it once more.
\end{tcolorbox}
\section[\texorpdfstring{\ru{капризный} (kapríznyy)}{капризный (kapríznyy)}]{\ru{капризный} (kapríznyy) — fussy, capricious}

\label{lesson:RU-C478-kapriznyy}



\begin{quote}
Before the new one: say the Russian for stubborn, then the Russian for stubbornness.
\end{quote}

\subsection*{You'll want to know: \ru{капризный}}
\textbf{\ru{капризный}} (\emph{kapríznyy}) — \textquotedblleft{}fussy, capricious\textquotedblright{}.

\begin{quote}
\ru{Капризный} \ru{ребёнок} \ru{плачет} \ru{из}-\ru{за} \ru{каждой} \ru{мелочи}.
\end{quote}

\begin{quote}
\emph{Kapríznyy rebyónok pláchet iz-za kázhdoy mélochi.}
\end{quote}

\begin{quote}
The fussy child cries over every little thing.
\end{quote}

Built on \emph{\ru{каприз}} — a whim.

\begin{cousinweb}
\emph{\ru{капризный}} is built on \emph{\ru{каприз}} 'whim', which was borrowed from French \emph{caprice}, itself from Italian \emph{capriccio} (EtymDB).
\end{cousinweb}

\begin{grammarlens}[title={What you've built}]
One more word for telling, explaining and describing at B1.
\end{grammarlens}

\subsection*{Your turn}
\begin{itemize}
\raggedright
  \item \emph{Say it:} \emph{kapríznyy}
  \item \emph{Say it:} \emph{kapríznyy}, once more
  \item \emph{Say it:} a sentence of your own with \emph{kapríznyy}
  \item \emph{Recall:} say the Russian for stubborn, then the Russian for stubbornness, then say \emph{kapríznyy} again
\end{itemize}

\begin{tcolorbox}[breakable,colback=teal!4,colframe=teal!35!black,title={Before you move on}]
What does \ru{капризный} mean? (\textquotedblleft{}Fussy, capricious\textquotedblright{}.) Say it once more.
\end{tcolorbox}
\section[\texorpdfstring{\ru{капризничать} (kapríznichat)}{капризничать (kapríznichat)}]{\ru{капризничать} (kapríznichat) — to play up (of a child)}

\label{lesson:RU-C478-kapriznichat}



\begin{quote}
Before the new one: say the Russian for stubbornness, then the Russian for fussy.
\end{quote}

\subsection*{You'll want to know: \ru{капризничать}}
\textbf{\ru{капризничать}} (\emph{kapríznichat}) — \textquotedblleft{}to play up (of a child)\textquotedblright{}.

\begin{quote}
\ru{Вечером} \ru{дочь} \ru{капризничает}, \ru{потому} \ru{что} \ru{устала}.
\end{quote}

\begin{quote}
\emph{Vécherom doch kapríznichayet, potomú chto ustála.}
\end{quote}

\begin{quote}
In the evening my daughter plays up because she's tired.
\end{quote}

Imperfective, colloquial; present \ru{капризничаю}.

\begin{cousinweb}
\emph{\ru{капризничать}} is built on \emph{\ru{каприз}} 'whim', which was borrowed from French \emph{caprice}, from Italian \emph{capriccio} (EtymDB).
\end{cousinweb}

\begin{grammarlens}[title={What you've built}]
One more word for telling, explaining and describing at B1.
\end{grammarlens}

\subsection*{Your turn}
\begin{itemize}
\raggedright
  \item \emph{Say it:} \emph{kapríznichat}
  \item \emph{Say it:} \emph{kapríznichat}, once more
  \item \emph{Say it:} a sentence of your own with \emph{kapríznichat}
  \item \emph{Recall:} say the Russian for stubbornness, then the Russian for fussy, then say \emph{kapríznichat} again
\end{itemize}

\begin{tcolorbox}[breakable,colback=teal!4,colframe=teal!35!black,title={Before you move on}]
What does \ru{капризничать} mean? (\textquotedblleft{}To play up (of a child)\textquotedblright{}.) Say it once more.
\end{tcolorbox}
\section[\texorpdfstring{\ru{ворчать} (vorchát)}{ворчать (vorchát)}]{\ru{ворчать} (vorchát) — to grumble}

\label{lesson:RU-C478-vorchat}



\begin{quote}
Before the new one: say the Russian for fussy, then the Russian for to play up.
\end{quote}

\subsection*{You'll want to know: \ru{ворчать}}
\textbf{\ru{ворчать}} (\emph{vorchát}) — \textquotedblleft{}to grumble\textquotedblright{}.

\begin{quote}
\ru{Дед} \ru{ворчит}, \ru{что} \ru{молодёжь} \ru{сидит} \ru{в} \ru{телефонах}.
\end{quote}

\begin{quote}
\emph{Ded vorchít, chto molodyózh sidít v telefónakh.}
\end{quote}

\begin{quote}
Grandad grumbles that young people are glued to their phones.
\end{quote}

Imperfective; present \ru{ворчу}, \ru{ворчишь}. Perfective \ru{проворчать}.

\begin{cousinweb}
EtymDB lists \emph{\ru{ворчать}}, glossed 'grumble', without a source; it is usually taken as imitative, like \emph{\ru{урчать}} 'to rumble'.
\end{cousinweb}

\begin{grammarlens}[title={What you've built}]
One more word for telling, explaining and describing at B1.
\end{grammarlens}

\subsection*{Your turn}
\begin{itemize}
\raggedright
  \item \emph{Say it:} \emph{vorchát}
  \item \emph{Say it:} \emph{vorchát}, once more
  \item \emph{Say it:} a sentence of your own with \emph{vorchát}
  \item \emph{Recall:} say the Russian for fussy, then the Russian for to play up, then say \emph{vorchát} again
\end{itemize}

\begin{tcolorbox}[breakable,colback=teal!4,colframe=teal!35!black,title={Before you move on}]
What does \ru{ворчать} mean? (\textquotedblleft{}To grumble\textquotedblright{}.) Say it once more.
\end{tcolorbox}
